%% The line mask, reported in two layers: the frequency coincidence, which is
%% a velocity criterion and nothing else, and the one species we do not read
%% as emission, which is an argument and is set out as one.
%% Owner: r16-mask.  Carries \label{app:linemask} (cited by 04_method.tex and
%% app_C_spatialnorm.tex) and \label{tab:masklines}.
%%
%% ★ THIS APPENDIX IS THE PAPER'S ONLY DEFINITION OF THE ATTRIBUTION RULE.
%% It owns the species list, the half-width and the FRAME.  The rule is the
%% star's own rest frame; the local standard of rest is evaluated as a check
%% and changes no disposition (masklsr_v520.py, L10).  Sec. 5.3 states the
%% rule in one clause and Appendix D points here; nothing else in the paper
%% may name a frame of its own, or there are two rules again.
\section{The line mask}
\label{app:linemask}

A crossing that falls on a known molecular transition has an astrophysical
explanation to hand, and the mask says which crossings those are.

\emph{The list.} The mask holds every transition of the \FrNSpec{} species
these setups cover -- the three CO isotopologues, HCN, HCO$^{+}$, CS, CN, SiO,
H$_2$CO, [C\,\textsc{i}] and hydrogen recombination -- inside the \McNIsland{}
searched intervals below \McEuMax\,K in upper-state energy, \FrNTrans{}
transitions (Table~\ref{tab:masklines}), main hyperfine components only and no
strength cut, from one Splatalogue query \citep{Splatalogue2007} of CDMS and
JPL \citep{Pickett1998,CDMS2005} dated \MkListDate.

\emph{The species are the ones the observers went to the telescope for, and
that is a limitation}: \TnCenPct{} per cent of the \NWinA{} Class~A windows
are centred within $\pm\MaskHalfKms$\,km\,s$^{-1}$ of a listed transition,
against \TnCtrlCenLo--\TnCtrlCenHi{} per cent for the same windows displaced,
so \emph{unattributed} means unattributed against the observers' own list. A
wider list -- every transition below \XlEuMax\,K of all \XlNSpecies{} species
with an interstellar or circumstellar detection, \XlNTrans{} in all -- covers
\XlBandPct{} per cent of the Class~A union against \MkMaskAPct{} per cent for
this one and attributes all \MkUnattr{} unattributed crossings with
\XlShareWidePct{} per cent of the control exceedances, emptying the category
for star and control alike; at that density proximity is not an
identification.

\emph{The pair that comes closest to an identification.} The \ZmNCross{}
unattributed crossings of $\beta$~Pictoris lie \ZmDvLo{} and
\ZmDvHi\,km\,s$^{-1}$ in the star's frame from methanol $5_K\to4_K$
components, offsets agreeing to \ZmDvAgree\,km\,s$^{-1}$ (\ZmChancePct{} per
cent by chance). Three spectroscopic measurements decide against them: the
brightest member, \ZmBright, would have been the window maximum and is not;
the weaker component is at \ZmRatioObs{} of the stronger where no temperature
allows more than \ZmRatioMax; and both drifts lie \ZmStepsOne{} and
\ZmStepsTwo{} grid steps from the zero a line at rest requires. Whether the
two could be one carrier is argued in Appendix~\ref{sec:pairacct}. The full
treatment is in \texttt{documentation/\allowbreak appendix\_C\_line\_mask.md} of the data
archive \citep{SIdeposit}.

\emph{Attribution is decided in the star's own rest frame, and the sign is
not the usual one.} A crossing's frequency is carried from the topocentric
value to the barycentre using its own block's pointing and mid-time, and then
to the stellar rest frame using the star's systemic velocity. That frame, and
no other, is where a coincidence is judged, a tolerance applied in the
observed frame spending most of its budget on the Earth's motion. The
\FrEpochNEpoch{} \FrEpochStar{} epochs show the transform working in the only
way that can be checked: \FrEpochTopoKms\,km\,s$^{-1}$ apart as observed, they
agree to \FrEpochStelKms\,km\,s$^{-1}$ in the star's frame, a twentieth of one
of their channels. Three are census crossings and the fourth is a hold-out
block. The observatory set all four executions to one source-frame frequency,
so the agreement bounds the part of the correction that changes with epoch,
which is the large part, and not a constant error in the systemic velocity.
Throughout, $\Delta v_\star = c\,(\nu_\star-\nu_0)/\nu_0$, so a positive
offset means the crossing lies \emph{above} the transition
in frequency, the opposite sign to the velocity convention of a spectral line.

\emph{Every crossing is also read in the local standard of rest, where
interstellar gas is at rest.} Toward \LsCloudStars{} the CO($1\to0$) survey
of \citet{DameThaddeus2022,Dame2001} shows emission at the crossing's own LSR
velocity (within \LsCloudDvMax\,km\,s$^{-1}$ of its peak, $\ge$\LsCloudTMin\,K),
so \LsNCloud{} of the \MkNCoinc{} coincidences are interstellar; elsewhere the
Planck CO map \citep{PlanckCO2016} stays below \LsPlanckMax\,K\,km\,s$^{-1}$.
HD~48370 is where the two readings part: its coincidence is
\LsHdDvStar\,km\,s$^{-1}$ from the star's own CO($2\to1$) but at
\LsHdVfeat\,km\,s$^{-1}$ in the LSR on a cloud line peaking at \LsHdVco, so
its extended emission is a cloud and not a disc.

Table~\ref{tab:masklines} gives the
mask species by species; the hydrogen recombination lines among them carry
no tabulated upper-state energy, and are blank in that column.

\begin{table}
\centering
\footnotesize
%% ★ v4.26 (r25-captions): 4 printed lines -> the author's 3.  The hydrogen
%% recombination lines' missing $E_{\rm u}$ MOVED to the paragraph above: it
%% explains a blank cell, which is a fact about the catalogue and not a
%% definition of a column.
\caption{The line mask, by species: the number $N$ of transitions, the
frequency range they span inside the searched intervals and their
upper-state energies.}
\label{tab:masklines}
%% GENERATED by maskso_v414.py -- do not hand-edit.
\begin{tabular}{@{}lrr@{\,}c@{\,}rr@{}}
\hline
Species & $N$ & \multicolumn{3}{c}{$\nu_0$ (GHz)} & $E_{\rm u}$ (K) \\
\hline
H & 5 & 92.03 & -- & 231.90 & -- \\
H$_2$CO & 13 & 101.33 & -- & 491.97 & 10--106 \\
C$^{18}$O & 3 & 109.78 & -- & 329.33 & 5--32 \\
CN & 50 & 113.12 & -- & 680.30 & 5--114 \\
CO & 3 & 115.27 & -- & 345.80 & 6--33 \\
CS & 3 & 195.95 & -- & 342.88 & 24--66 \\
SiO & 2 & 217.10 & -- & 347.33 & 31--75 \\
$^{13}$CO & 8 & 220.40 & -- & 330.59 & 16--32 \\
HCN & 7 & 354.50 & -- & 354.51 & 43--43 \\
HCO$^{+}$ & 1 & 356.73 & -- & 356.73 & 43 \\
{[C\,\textsc{i}]} & 1 & 492.16 & -- & 492.16 & 24 \\
\hline
\end{tabular}%

\end{table}

\emph{The coincidences, and the one species left out.} \MkNCoinc{} of the
\NCross{} crossings fall within the mask in their own star's frame, and these
are the \MkAttr{} attributions. Sulphur monoxide is deliberately not listed --
no debris disc has a reported SO detection, so masking its \SoNTrans{}
transitions would have declared no constraint over bandwidth the survey does
constrain -- and the \SoNCrossWord{} crossings within the half-width of one
(\SoStars, at \SoDvList\,km\,s$^{-1}$) are reported as unattributed,
\SoNBelowScreen{} below the rank screen. Nor do the tests made find
interstellar SO: in the LSR the offsets are \SoDvLsrList\,km\,s$^{-1}$, none within
the cloud test, and no CO($3\to2$) feature above the trigger accompanies any.
The hold-out is treated identically, \HoNAttr{} attributed and \HoNUnattr{}
unattributed (\S\ref{sec:holdoutsearch}).

\emph{The half-width, and what it costs.} Keplerian motion where these
molecules survive, $\lesssim$13\,km\,s$^{-1}$ beyond 10\,au, sets it at
$\pm\MaskHalfKms$\,km\,s$^{-1}$. \emph{The cost is stated on the union of
searched sky frequencies}, the axis on which coverage is quoted throughout:
the frequencies within that half-width of a listed transition's rest
frequency, merged where transitions overlap, are \MkMaskAGHz\,GHz of the
\CvUnionA\,GHz Class~A union, \MkMaskAPct{} per cent, and \FrMaskLostGHz\,GHz
of \FrMaskUnionGHz\,GHz over both classes, \FrMaskPct{} per cent; the wider
list's share is measured the same way. A union over windows toward different
stars has no single rest frame, so this describes the coverage and not any one
star; a limit is formed instead over its own window's channels outside the
mask in that star's own frame (\S\ref{sec:summary}). Attributing on either
frame would mask \LsMaskAddPct{} per cent more for no further attribution.

Of the \MkAttrExcl{} labelled crossings not recovered (\S\ref{sec:etacrv}),
one rests on a velocity instead, its sightline carrying interstellar emission
\LsLupDv\,km\,s$^{-1}$ from the crossing, with \RqAjNgeLo--\RqAjNgeHi{} of
the \NCtrl{} control positions above the star in every repeat block.
